% ICML 2027 submission — EDA-Flow-Bench (cognition paper, v1.7 rewrite)
% 双盲投稿：使用 [submitted] 选项（匿名）
% camera-ready 时改为 [accepted] 并填作者信息
%
% 编译：pdflatex main && bibtex main && pdflatex main && pdflatex main
%
% v1.7 重写：从"卖 bench"重定位为"认知命题 + 情绪机制 + 可证伪假设"
% 正文 ≤ 6 页（不含 references + appendix）
%
% 【中文】main.tex 覆盖说明（2026-11-18 中文对照令）：
% 本文件无独立英文正文段落——仅含 \documentclass/宏定义/标题/作者块与 \input 章节指令。
% 按令「标题/作者/宏定义旁注不必译」，故本文件不插入中文对照；中英对照落在各 sections/*.tex 与
% appendix/*.tex 内。标题英文原样保留（双盲投稿标题不译）。

\documentclass{article}

% ICML 2027 Author Kit 尚未发布；当前使用仓库自带的 icml2026.sty/.bst（接口一致）。
\usepackage{hyperref}
\usepackage{icml2026}

% 基础包
\usepackage{amsmath,amssymb}
\usepackage{graphicx}
\usepackage{booktabs}
\usepackage{algorithm}
\usepackage{algorithmic}
\usepackage{xspace}
\usepackage{multirow}
\usepackage{tabularx}
\usepackage{enumitem}

% 排版微调（2026-11-19 用户令 · Overfull/Underfull 修复）：
%   \emergencystretch 让段落断行在难断处最多额外拉伸 3em 而不报 badness，
%   显著减少 Underfull \hbox（段落级）；\raggedbottom 允许双栏底部不强制对齐，
%   消除 "Underfull \vbox (badness 10000) while \output is active"（栏底填充警告）。
\setlength{\emergencystretch}{4em}
\raggedbottom
% ragged X 列：用于 tabularx 中含不可断数学串（如 M7$<$M2$-$2pp）的单元格，
% 避免单 token 行被强制拉伸而报 Underfull \hbox (badness 10000)。
\newcolumntype{Y}{>{\raggedright\arraybackslash}X}

% 宏定义
\newcommand{\bench}{EDA-Flow-Bench\xspace}
\newcommand{\mode}{design\_flow\xspace}
\newcommand{\eg}{e.g.,\xspace}
\newcommand{\ie}{i.e.,\xspace}
\newcommand{\etal}{\textit{et al.}\xspace}
\newcommand{\TBD}{\textbf{TBD}\xspace}

% 评测集（2026-11-23 用户令：单一 500 任务集，不出现 Full/Lite 两层）：
%   500 任务，覆盖 13 个 EDA domain，按链长/工具链深度分层。
\newcommand{\Ncases}{500}
% M0 基线 Pass@1（DeepSeek-V4-Pro-cloud · 3-run mean ± std · 2026-11-24）
%   r1=21.78%(105/482) · r2=19.75%(96/486) · r3=21.02%(103/490) · mean=20.85% ± 1.03%
\newcommand{\MzeroPass}{20.85\%}
\newcommand{\MzeroStd}{1.03}
% 4 骨干对比（Lite 500 · Pass@1 · 3-run mean ± std）
\newcommand{\DSvPro}{20.85\%}
\newcommand{\DSvProStd}{1.03}
\newcommand{\DoubaoPro}{10.18\%}
\newcommand{\DoubaoProStd}{1.18}
\newcommand{\DoubaoLite}{4.68\%}
\newcommand{\DoubaoLiteStd}{0.77}
\newcommand{\DoubaoMini}{2.39\%}
\newcommand{\DoubaoMiniStd}{0.19}

% 图表路径
\graphicspath{{figures/}}

% ===== 标题与作者 =====
\begin{document}

\twocolumn[
  \icmltitle{Affective Gating of Long-Term Memory in Long-Chain EDA Agents: A Falsifiable Study}

  \icmltitlerunning{Affective Memory for EDA Flow Agents}

  % 双盲：不写真实作者
  \begin{icmlauthorlist}
    \icmlauthor{Anonymous Submission}{anon}
  \end{icmlauthorlist}

  \icmlaffiliation{anon}{Anonymous Institution}

  \icmlcorrespondingauthor{Anonymous Submission}{anonymous@example.org}

  \icmlkeywords{EDA, LLM agents, long-chain orchestration, affective memory, falsifiable hypotheses}

  \vskip 0.3in
]
\printAffiliationsAndNotice{}

% ===== 主体章节（按 v1.6 计划书 §6.3 结构）=====
\begin{abstract}
Frontier LLM agents collapse on long-chain EDA design-flow orchestration: on \bench, a single frontier backbone reaches only \textbf{\MzeroPass{} Pass@1} (62/620 decided runs, frozen 696-run snapshot) versus \textbf{\ScriptPass{}} on single-script EDA benchmarks---a 7.6$\times$ drop. We hypothesize that the gap is not merely a model-capacity ceiling but a \emph{memory and self-improvement} deficit that is mechanistically tractable. We pre-register \textbf{3 core falsifiable hypotheses} (H1: governed long-term memory + recursive self-improvement yields transferable gains on long-chain flow; H2: the gains arise from three affective-memory-inspired computational mechanisms---surprise-gated writing, weapon-focus compression, state-aware retrieval---rather than mere prompt lengthening; H3: the prior finding that single-agent $\geq$ multi-agent replicates in the flow setting) decomposed into \textbf{7 sub-hypotheses}, and design a \textbf{13-arm ablation} including an equal-token killer control (N1) and a cross-domain transfer contamination probe (M4). The three mechanisms map cognitive-neuroscience principles (amygdala--hippocampus arousal consolidation, weapon-focus central/peripheral trade-off, mood-congruent retrieval) to computable formulas: $w_{\text{write}}=\alpha\cdot\text{surprise}+\beta\cdot\text{penalty}$, $r_{\text{cp}}=|\text{central}|/|\text{total}|$, $\text{score}=\text{sim}+\gamma\cdot\text{state\_match}$. Results for the M0 baseline are available from the frozen snapshot; all other arms are to be evaluated after the pre-registration date (\PreregDate). Negative results will be reported as publishable boundary conditions on RSI and affective-memory mechanisms in symbolic tool orchestration. We release a protocol white-paper (harness protocol + criteria + desensitized task specs) under a controlled NDA review package, without releasing executable Lite-set assets.
\end{abstract}

\section{Introduction}
\label{sec:intro}
% 【中文】\section{引言}

LLM agents have mastered single-function EDA code generation, yet they collapse on \emph{long-chain design-flow orchestration}: multi-step tasks that require driving live commercial EDA tools (schematic capture, SPICE simulation, DRC/LVS, synthesis) through dozens of tool calls, inspecting outputs, and recovering from errors. On \bench, a single frontier backbone (DeepSeek-V4-Pro) reaches only \MzeroPass{}$\pm$\MzeroStd{} Pass@1 over 3 runs on the \Ncases{}-task set, while three weaker backbones range from \DoubaoMini{} to \DoubaoPro{}---a 2.0$\times$ spread non-uniform across domains, and the hardest domains are precisely those demanding the deepest live-tool interaction.
% 【中文】
% LLM 智能体已掌握单功能 EDA 代码生成，但它们在\emph{长链设计流编排}上崩溃：这类多步任务要求驱动
% 真实商用 EDA 工具（原理图录入、SPICE 仿真、DRC/LVS、综合）经数十次工具调用、检视输出并从错误中
% 恢复。在 \bench 上，单一前沿骨干（DeepSeek-V4-Pro）在 \Ncases{}-任务集上 3 轮仅达到
% \MzeroPass{}$\pm$\MzeroStd{} Pass@1，三个较弱骨干从 \DoubaoMini{} 到 \DoubaoPro{}
% 不等——2.0$\times$ 跨域不均匀差距，而最难的领域恰恰是那些要求最深实时工具交互的领域。

We argue this is not merely a model-capacity problem but a \emph{memory and self-improvement} deficit. Long-chain flow tasks exceed any context window: an agent that cannot consolidate lessons across sessions---which errors recur, which tool sequences work---is condemned to re-discover the same knowledge every run. Recursive self-improvement (RSI)---expanding a skill library and refining the harness from experience---is only meaningful if memory persists to accumulate that experience.
% 【中文】
% 我们论证这并非单纯的模型能力问题，而是一个\emph{记忆与自我改进}缺失。长链流程任务超出任何上下文
% 窗口：一个无法跨会话巩固经验教训（哪些错误反复出现、哪些工具序列奏效、哪些项目约定要紧）的智能体，
% 注定每次运行都要重新发现同样的知识。递归自我改进（RSI）——扩展技能库并从经验中改进执行框架本身的
% 能力——只有在记忆跨会话持久存在以累积经验时才有意义。

\textbf{Contributions.} (1) We \emph{introduce} \bench{}, a benchmark of \Ncases{} long-chain EDA design-flow tasks across 13 domains, each derived from real IC-design engineer coding-agent traces in a production environment, with a \mode{} contract (submit.zip + $\geq$8 executable assertions + a PreToolUse anti-leakage hook) and an M0 baseline \MzeroPass{}$\pm$\MzeroStd{} Pass@1 (3-run, DeepSeek-V4-Pro); no data or code is released, with reproducibility supported via the protocol in this paper and an NDA-gated review package (\S\ref{sec:benchmark}). (2) We \emph{formulate} 3 core falsifiable hypotheses (H1 functional / H2 mechanistic / H3 boundary) decomposed into 7 sub-hypotheses, each with falsification conditions, and formalize three cognitive-neuroscience mechanisms---arousal $\rightarrow$ surprise/penalty gating, weapon-focus $\rightarrow$ central-peripheral compression, mood-congruence $\rightarrow$ state-aware retrieval---as computable mappings (\S\ref{sec:hypotheses}, \S\ref{sec:harness}). (3) We \emph{design} a 10-arm main ablation (3 appendix arms) with an equal-token killer control (N1, excluding ``prompt just got longer'') and a cross-domain transfer contamination probe (M4, with a machine-judgeable leak-rate audit); negative results are reported as publishable boundary conditions.
% 【中文】
% \textbf{贡献。}(1) 我们\emph{引入}\bench{}——一个含 \Ncases{} 个长链 EDA 设计流任务、跨 13 个领域的
% 基准，每个任务源自真实 IC 设计工程师在生产环境中的 coding-agent trace，含 \mode{} 契约
%（submit.zip + $\geq$8 个可执行断言 + PreToolUse 反泄漏钩子）与 M0 基线
% \MzeroPass{}$\pm$\MzeroStd{} Pass@1（3-run，DeepSeek-V4-Pro）；不发布数据/代码，复现性由本论文协议
% 描述 + NDA 受控评审包支持（\S\ref{sec:benchmark}）。(2) 我们\emph{提出}
% 3 条核心可证伪假设（H1 功能性 / H2 机制性 / H3 边界性）下设 7 条子假设、各含证伪条件，
% 并形式化三个认知神经科学机制——情绪觉醒 $\rightarrow$ 惊异/惩罚门控、武器焦点 $\rightarrow$
% 中心-外围压缩、情绪一致性 $\rightarrow$ 状态感知检索——为可计算映射
%（\S\ref{sec:hypotheses}、\S\ref{sec:harness}）。(3) 我们\emph{设计} 10 臂主消融（3 臂附录），
% 含等 token 杀手对照（N1，排除"提示变长"替代解释）与跨域迁移污染探针（M4，含机器可判的
% 答案泄露率审计）；阴性结果作为可发表的边界条件报告。

\textbf{Experimental status.} The M0 baseline is evaluated (\MzeroPass{}$\pm$\MzeroStd{} Pass@1, 3-run, DeepSeek-V4-Pro); the other arms (M1--M8, N1--N2) are to be evaluated.
% 【中文】
% \textbf{实验状态。} M0 基线已评测（\MzeroPass{}$\pm$\MzeroStd{} Pass@1，3-run，DeepSeek-V4-Pro）；
% 其他臂（M1--M8、N1--N2）待评测。

\section{Background: From Script to Flow}
\label{sec:background}

\textbf{The grounding gap.} EDA scripting benchmarks (VerilogEval, RTLLM) evaluate isolated code generation: given a spec, produce Verilog. The agent makes 1--3 tool calls, never inspects a live tool's output, and never recovers from an error. \bench{} introduces \emph{design-flow orchestration}: the agent must produce a deliverable (\texttt{submit.zip}) by orchestrating remote commercial EDA tools through an MCP bridge, scored against $\geq 8$ executable assertions. The median run issues 26.5 tool calls (IQR 13.8--38.0) across 20.5 iterations, with 41.2\% of calls going through the MCP bridge to live EDA tools (47 instrumented pilot traces; \S\ref{sec:analysis}). This is qualitatively harder: the agent must ground its actions in tool outputs it has never seen during training, chain dependent steps, and recover from failures mid-flow.

\textbf{Why memory and RSI may matter more in flow.} On single-script tasks, in-context learning within one trace often suffices---the entire problem fits in the window. On long-chain flow, the relevant knowledge (tool quirks, project conventions, failure patterns) accumulates \emph{across} sessions and exceeds any window. This makes persistent memory a potential binding constraint: an agent that can consolidate ``which tool sequences work'' and ``which errors recur'' into a cross-session memory may close the grounding gap, while one that re-derives everything from scratch each run is bounded. Recursive self-improvement (RSI)---expanding a skill library and refining the harness from experience---is only meaningful if memory persists to accumulate that experience. Our prior work on script-layer EDA (DATE27, \cite{PLACEHOLDER_DATE27_ZHULONG}; ZhuLong-DAC, \cite{liu2026runreadbackrefine}) showed memory + RSI gains on short-chain tasks; whether these \emph{transfer} to long-chain flow is the central question of this paper.

\textbf{From biological affect to computational mapping.} Cognitive neuroscience establishes that emotional arousal modulates memory consolidation (amygdala--hippocampus interaction; \cite{mcgaugh2004}), that high arousal triggers a central/peripheral encoding trade-off (weapon-focus; \cite{loftus1987}), and that retrieval is mood-congruent and state-dependent \cite{bower1981}. We map these three principles to computable mechanisms for an EDA agent's memory architecture (\S\ref{sec:harness}), turning ``does memory help?'' into ``\emph{which} mechanism brings the gain, and under what conditions does it vanish?''

\section{Task Modality and Benchmark (Carrier)}
\label{sec:benchmark}
% 【中文】\section{任务模态与基准（载体）}

\bench{} is a benchmark of \textbf{\Ntasks{} tasks} across \textbf{13 EDA domains}---from schematic/netlist to SPICE simulation, physical verification, and digital synthesis---each derived from \emph{real coding-agent traces} of IC design engineers in a production environment. We use \bench{} as a \emph{carrier} for falsifiable cognition claims, not as the contribution itself.
% 【中文】
% \bench{} 是一个含 \textbf{\Ntasks{} 个任务}、跨越 \textbf{13 个 EDA 领域}的基准——从原理图/网表到
% SPICE 仿真、物理验证与数字综合——每个任务源自生产环境中 IC 设计工程师的\emph{真实编程智能体轨迹}。
% 我们将 \bench{} 用作可证伪认知主张的\emph{载体}，而非贡献本身。

\textbf{Design-flow contract.} Each task requires the agent to produce \texttt{submit.zip} by orchestrating remote EDA tools through an MCP bridge, scored against $\geq 8$ executable assertions with a \texttt{PreToolUse} hook that prevents answer leakage (reference/ and check.py never enter the workdir). The contract deliberately does not specify language or host---the agent must discover the right tool chain. Detailed benchmark construction is in Appendix~\ref{sec:benchmark-details}.
% 【中文】
% \textbf{设计流契约。} 每个任务要求智能体通过 MCP 网桥编排远程 EDA 工具产出 \texttt{submit.zip}，
% 按 $\geq 8$ 条可执行断言评分，并用 \texttt{PreToolUse} 钩子防止答案泄漏（reference/ 与 check.py
% 永不进入工作目录）。该契约故意不规定语言或主机——智能体必须自行发现正确的工具链。详细基准构建见
% 附录~\ref{sec:benchmark-details}。

\textbf{Evaluation subset.} Experiments use a Full-set subset (top-120 by grounding depth / chain length, internally evaluated, not released). A frozen 696-run snapshot (2026-11-09; DeepSeek-V4-Pro, Cline v3.0.51, \texttt{thinking=low}, 2400s timeout, 8 workers) provides the M0 baseline anchor: \MzeroPass{} Pass@1 (62 PASS / 558 FAIL / 70 ERROR / 6 ABSENT).
% 【中文】
% \textbf{评测子集。} 实验使用 Full 集子集（按接地深度/链长取 top-120，内部评测，不发布）。一份冻结
% 的 696-run 快照（2026-11-09；DeepSeek-V4-Pro、Cline v3.0.51、\texttt{thinking=low}、2400s 超时、
% 8 workers）提供 M0 基线锚点：\MzeroPass{} Pass@1（62 PASS / 558 FAIL / 70 ERROR / 6 ABSENT）。

\textbf{Open-source boundary (v1.4 narrowing).} We release a \emph{protocol white-paper} (harness protocol + criteria + desensitized task specifications) and an \emph{NDA-controlled review package} for reviewers, but do \textbf{not} release executable Lite-set assets or a mock-EDA mode. The falsifiability of our claims does not depend on releasing proprietary EDA tool licenses: the protocol and criteria are sufficient for an independent party to re-derive an equivalent benchmark on open-source EDA tools (Yosys/OpenROAD/Magic) and test the same hypotheses (\S\ref{sec:reproducibility}).
% 【中文】
% \textbf{开源边界（v1.4 收缩）。} 我们发布\emph{协议白皮书}（执行框架协议 + 判据 + 脱敏任务规格）与
% \emph{NDA 受控评审包}供评审者使用，但\textbf{不}发布可执行的 Lite 集资产或 mock-EDA 模式。我们主张的
% 可证伪性不依赖于发布专有 EDA 工具许可：协议与判据足以让独立方在开源 EDA 工具（Yosys/OpenROAD/Magic）
% 上重新推导出等价基准并检验同样假设（\S\ref{sec:reproducibility}）。

\section{Harness with Memory and Self-Improvement}
\label{sec:harness}

\subsection{Three-Layer Memory (L0/L1/L2)}
\label{sec:memory-arch}

We extend the three-layer memory architecture from our prior script-layer work \cite{liu2026runreadbackrefine} to the flow setting. \textbf{L0} (working memory): the current trace's context window. \textbf{L1} (session memory): distilled lessons from the current run (error $\rightarrow$ fix mappings, tool-sequence templates), with reliability gating---ground-truth anchoring (a binary validator must cover self-reported success), disconfirmation filtering (2 refutations $\rightarrow$ evict), 5-round unconfirmed demotion, and held-out gated evolution (regress $>5$pp $\rightarrow$ rollback). \textbf{L2} (cross-session skill library): durable, versioned skills indexed by abstraction/reusability/actionability/faithfulness, with anti-leakage retrieval (names-only injection, docstring generalization, problem-side embedding) to prevent cross-task answer copying.

\subsection{Affective-Memory Mechanisms: From Biology to Computation}
\label{sec:affective}

The L1/L2 architecture above is \emph{functional}: it asks ``does memory help?'' but not ``\emph{which} computational mechanism brings the gain?'' We introduce three mechanisms inspired by cognitive-neuroscience principles of emotional memory, each giving a \emph{computable} mapping that can be ablated with a mechanism-level negative control (equal-token budget, mechanism off).

\textbf{Mechanism 1: Surprise/Penalty-Gated Writing.} Humans do not log every event uniformly; emotionally arousing events---unexpected failures or breakthroughs---are preferentially consolidated via amygdala--hippocampus interaction \cite{mcgaugh2004}. We map this to a \emph{write-gating weight}:
\begin{equation}
\label{eq:write-gate}
w_{\text{write}}(e) = \alpha \cdot \text{surprise}(e) + \beta \cdot \text{penalty}(e),
\end{equation}
where $\text{surprise}(e) = \text{KL}(p_{\text{posterior}} \| p_{\text{prior}})$ is the Bayesian surprise of event $e$ relative to the memory bank's prior \cite{ittibaldi2009}, and $\text{penalty}(e)$ is the sunk cost (tool calls / tokens consumed before a mid-chain crash). The agent writes preferentially to high-$w_{\text{write}}$ events; $\alpha,\beta$ are calibrated by ablation (M5 vs.\ M2 uniform sampling), with total memory budget held equal to control for ``just remembering more.''

\textbf{Mechanism 2: Weapon-Focus Compression.} Under high arousal, the brain consolidates \emph{central} details (the threatening stimulus) at the expense of \emph{peripheral} details \cite{loftus1987}. We map this to a compression ratio:
\begin{equation}
\label{eq:weapon-focus}
r_{\text{cp}} = \frac{|\text{central tokens}|}{|\text{total tokens}|},
\end{equation}
where \emph{central} = error stack trace + root-cause causal chain, and \emph{peripheral} = intermediate tool-call logs. When surprise exceeds a threshold (same trigger as M5), the agent writes only central content into L1/L2, stripping peripheral noise. This predicts: smaller memory $\wedge$ equal-or-better Pass@1.

\textbf{Mechanism 3: State-Aware Retrieval.} Mood-congruent and state-dependent recall \cite{bower1981} predict that an agent in a ``failing streak'' state should preferentially retrieve defensive pitfall-avoidance lessons, while a ``smooth progress'' state should retrieve success templates. We map this to a retrieval score:
\begin{equation}
\label{eq:state-retrieval}
\text{score}(q, m) = \text{sim}(q, m) + \gamma \cdot \text{state\_match}(\text{state}(q), \text{state}(m)),
\end{equation}
where $\text{state} = (\text{recent-}N\text{-step failure rate}, \text{current tool domain}, \text{chain-stage})$ and $\gamma$ is calibrated by ablation (M6 vs.\ M3 state-agnostic retrieval, same memory content).

Each mechanism has a \emph{mechanism-level negative control}: turn the mechanism off while holding the token budget equal, isolating ``the mechanism helps'' from ``more context helps.'' This is stricter than the functional-level control (M3 vs.\ M0) and directly addresses the reviewer objection that memory gains are just prompt-lengthening. Full mechanism derivations, biological-evidence citations, and the AI-affective-memory literature positioning (Generative Agents \cite{park2023}, EM-LLM, PsychoAgent) are in Appendix~\ref{sec:affective-appendix}.

\section{Falsifiable Hypotheses and Ablation Design}
\label{sec:hypotheses}
% 【中文】\section{可证伪假设与消融设计}

We formulate \textbf{3 core falsifiable hypotheses} (H1 functional / H2 mechanistic / H3 boundary) decomposed into \textbf{7 sub-hypotheses}, organized in a \emph{total--mechanism--boundary} three-layer structure, each with falsification conditions.
% 【中文】
% 我们提出 \textbf{3 条核心可证伪假设}（H1 功能性 / H2 机制性 / H3 边界性），分解为 \textbf{7 条子假设}，
% 按\emph{总-机制-边界}三层结构组织，每条含证伪条件。

\subsection{Three-Layer Structure}
% 【中文】\subsection{三层结构}
% 【中文】表 1：三层假设结构（列名依次：层级、核心假设、子假设）。tabularx+X 列防双栏越 margin。

\begin{center}
\small
\begin{tabularx}{\columnwidth}{@{}lXX@{}}
\toprule
\textbf{Layer} & \textbf{Core Hyp.} & \textbf{Sub-Hypotheses} \\
\midrule
L1 Total & H1: memory+RSI $\to$ transferable gain & H1a, H1b, H1c \\
L2 Mechanism & H2: gain from 3 affective mechanisms & H2a, H2b, H2c \\
L3 Boundary & H3: single-agent $\geq$ multi-agent & H3 (independent) \\
\bottomrule
\end{tabularx}
\end{center}

\textbf{H1 (total).} Governed long-term memory + RSI yields transferable Pass@1 gains on long-chain EDA flow. \emph{Falsification (disjunctive):} any of H1a/H1b/H1c falsified $\Rightarrow$ H1 falsified.
% 【中文】
% \textbf{H1（总）。} 受治理的长期记忆 + RSI 在长链 EDA 流程上带来可迁移的 Pass@1 增益。
% \emph{证伪（析取式）：} H1a/H1b/H1c 任一被证伪 $\Rightarrow$ H1 被证伪。

\textbf{H2 (mechanism).} H1's gain arises from the three affective-memory mechanisms (\S\ref{sec:affective}), not prompt lengthening (N1 excludes). \emph{Falsification (conjunctive):} all of H2a/H2b/H2c falsified $\Rightarrow$ H2 falsified; 1--2 surviving $\to$ ``partial mechanisms effective.''
% 【中文】
% \textbf{H2（机制）。} H1 的增益源于三个情绪记忆机制（\S\ref{sec:affective}），而非提示延长（N1 排除）。
% \emph{证伪（合取式）：} H2a/H2b/H2c 全部被证伪 $\Rightarrow$ H2 被证伪；1--2 个机制存活 $\to$
% "部分机制有效"。

\textbf{H3 (boundary, appendix).} The prior finding (single-agent $\geq$ multi-agent; autonomous-spawn $\geq$ user-custom) replicates in the flow setting. \emph{Independent} of H1/H2. Both outcomes publishable: replication $\to$ consolidates; non-replication $\to$ bounds the single-agent advantage to short chains.
% 【中文】
% \textbf{H3（边界，附录）。} 先前发现（单智能体 $\geq$ 多智能体；自主 spawn $\geq$ 用户定制）在流程场景
% 可复现。\emph{独立于} H1/H2。两种结果均可发表：复现 $\to$ 巩固；不可复现 $\to$ 将单智能体优势限定于
% 短链。

\subsection{Sub-Hypotheses (compact)}
% 【中文】\subsection{子假设（紧凑版）}
% 【中文】表 2：子假设紧凑表（列名依次：子假设、主张（一行）、证伪条件（一行））

\begin{center}
\scriptsize
\setlength{\tabcolsep}{3pt}
\begin{tabularx}{\columnwidth}{@{}lYY@{}}
\toprule
\textbf{Sub} & \textbf{Claim (one line)} & \textbf{Falsification (one line)} \\
\midrule
H1a & Memory lifts long-chain Pass@1 (M3 vs M0, $\geq$3pp) & M3$-$M0$<$1pp or McNemar $p>0.05$ or reversal \\
H1b & RSI gain grows monotonically with chain length ($\Delta_{\text{long}}\geq\Delta_{\text{mid}}\geq\Delta_{\text{short}}$, $\geq$3pp) & monotonicity broken or $\Delta_{\text{long}}-\Delta_{\text{short}}$ CI contains 0 \\
H1c & Cross-session long-term memory is necessary for cross-session RSI (no-memory cross-session RSI gain $\leq$0; interaction $\geq$2pp) & S1$>$M0+2pp or interaction CI contains 0 \\
\midrule
H2a & Surprise-gated writing $>$ uniform (M5 vs M2, $\geq$3pp, equal token) & M5$-$M2$<$1pp or reversal \\
H2b & Weapon-focus compression $>$ full write (M7 vs M2, $\geq$2pp, memory $\leq$0.6$\times$) & M7$<$M2$-$2pp or memory $\geq$0.8$\times$M2; {\em weak support}: memory 0.6--0.8$\times$M2 \& $|$M7$-$M2$|<$2pp \\
H2c & State-aware retrieval $>$ state-agnostic (M6 vs M3, $\geq$3pp on fail-streak subset) & fail-streak M6$-$M3$<$1pp or harm in smooth subset \\
\midrule
H3 & Single-agent $\geq$ autonomous-spawn in flow (nested in M3+S1) & autonomous-spawn $>$ single-agent +3pp (McNemar $p<0.05$) \\
\bottomrule
\end{tabularx}
\end{center}

\textbf{Statistics.} McNemar (paired, direction-consistent) + paired bootstrap 95\% CI (1000 resamples); $\alpha=0.05$ (Bonferroni $\alpha=0.05/3=0.0167$ for the 3-hypothesis family). Minimum matrix: 120 tasks/arm (top-120 of the Lite \LiteSet{}-task set by chain length); \emph{ideal matrix = Lite $\times$ 3 repeats} (\LiteSet{} tasks $\times$ 3, power $\geq$0.8). \emph{Power positioning:} at $n{=}120$, McNemar power for a 3pp effect is $\approx$0.62, so minimal-matrix results are reported as \textbf{pilot signal}, not primary verdicts; H1/H2 verdicts are deferred to the ideal matrix (Lite $\times$ 3, power $\geq$0.8). M6's fail-streak subset ($\sim$40 tasks at $n{=}120$, power $<$0.3) is exploratory; H2c is deferred to the ideal matrix (fail-streak $\geq$80). The abstract reports no specific pp numbers for arms.
% 【中文】
% \textbf{统计。} McNemar（配对、方向一致）+ 配对 bootstrap 95\% CI（1000 次重采样）；$\alpha=0.05$
%（3 假设族 Bonferroni $\alpha=0.05/3=0.0167$）。最小矩阵：120 任务/臂（按链长取 Lite \LiteSet{}-任务集
% 的 top-120）；\emph{理想矩阵 = Lite $\times$ 3 重复}（\LiteSet{} 任务 $\times$ 3、功效 $\geq$0.8）。
% \emph{功效定位：} $n{=}120$ 下 3pp 效应的 McNemar 功效 $\approx$0.62，
% 故最小矩阵结果作为\textbf{试点信号}报告，不作主要判据；H1/H2 主要判据延至理想矩阵
%（Lite $\times$ 3、功效 $\geq$0.8）。M6 的失败连续子集（$n{=}120$ 时 $\sim$40 题、3pp 功效 $<$0.3）
% 仅为探索性；H2c 判据延至理想矩阵（失败连续 $\geq$80 题）。摘要不报告各臂具体 pp 数字。

\subsection{10-Arm Main Ablation (3 Appendix)}
\label{sec:arms}
% 【中文】\subsection{10 臂主消融（3 臂附录）}
% 【中文】表 3：10 臂主消融设计（列名依次：臂、对应假设、变量、判据）。M2 作 H2a/H2b 对照臂（不进主表）；M1/M8 附录。

\begin{center}
\scriptsize
\setlength{\tabcolsep}{2.5pt}
\begin{tabularx}{\columnwidth}{@{}llYY@{}}
\toprule
\textbf{Arm} & \textbf{Hyp} & \textbf{Variable} & \textbf{Criteria} \\
\midrule
M0 & H1a & no memory (baseline) & anchor = M0 Pass@1 on same top-120-of-Lite subset (\MzeroAnchor{}, 12/119 decided, doubao-seed-2.0-pro-cloud); 696-run \MzeroPass{} is Full-set descriptive only \\
M3 & H1a & auto memory + reliability gating & \textbf{M3$\geq$M0+3pp} \\
S1 & H1c & RSI, no memory; inter-session persistence explicitly forbidden (incl.\ harness-level file residue) & S1$\leq$M0 (cross-session RSI null under no inter-session persistence) \\
M3+S1 & H1b,c & memory + RSI & interaction $\geq$2pp; $\Delta_{\text{long}}\geq\Delta_{\text{mid}}\geq\Delta_{\text{short}}$ \\
\midrule
N1 & H1a & killer: equal-token random padding & M3$\geq$N1+3pp (excludes ``prompt longer'') \\
N2 & H2$^-$ & error/irrelevant memory & N2$\leq$M0; M3$\geq$N2+5pp \\
M4 & H1+H2 & transfer: train-domain write, target-domain read & $\geq$M0+2pp; leak rate $=0$ (machine-judgeable proxy: n-gram overlap of memory entries vs target-domain prompt/filename/assertion literals; manual only borderline) \\
\midrule
M5 & H2a & surprise-gated writing ($w=\alpha\cdot$surprise$+\beta\cdot$penalty) & M5$\geq$M2+3pp; token $\Delta\leq$10\% \\
M6 & H2c & state-aware retrieval (score=sim$+\gamma\cdot$state\_match) & fail-streak M6$\geq$M3+3pp (minimal-matrix exploratory: $\sim$40 tasks, power $<$0.3; H2c verdict deferred to ideal matrix, fail-streak $\geq$80) \\
M7 & H2b & weapon-focus compression ($r_{\text{cp}}$) & M7$\geq$M2+2pp; memory $\leq$0.6$\times$ \\
\bottomrule
\end{tabularx}
\end{center}

\textbf{Appendix arms.} M2 (auto memory, no gating) is the H2a/H2b control arm---it appears in H2 criteria rows (M5 vs M2, M7 vs M2) but is \emph{not} a main-table arm; its prior H1a weak-control criteria (M2$\geq$M0+2pp and M3$\geq$M2+1pp) are \textbf{withdrawn} (1pp is below detectable effect at $n{=}120$). M1 (static SOP memory) and M8 (negative-valence bias) are appendix robustness checks.
% 【中文】
% \textbf{附录臂。} M2（自动记忆、无门控）作为 H2a/H2b 对照臂评测——出现在 H2 判据行（M5 vs M2、M7 vs M2），
% 但\emph{不}是主表臂；其先前的 H1a 弱对照判据（M2$\geq$M0+2pp 与 M3$\geq$M2+1pp）\textbf{撤回}
%（1pp 在 $n{=}120$ 下低于可检效应，§5.3）。M1（静态 SOP 记忆）与 M8（负效价偏倚）为附录稳健性检验，不进主 10 臂表。

All arms use the same \texttt{check.py}, same $\geq$8 assertions, same 2400s timeout, same backbone (\textbf{doubao-seed-2.0-pro-cloud}; the frozen 696-run snapshot used DeepSeek-V4-Pro, which has since been deprecated). M1--M8 and N1--N2 results are \TBD (pending evaluation); M0's minimal-matrix anchor is the same-subset M0 on the top-120-of-Lite (\MzeroAnchor{}, 12/119 decided), not the 696-run Full-set \MzeroPass{} (descriptive baseline snapshot).
% 【中文】
% 所有臂使用相同的 \texttt{check.py}、相同的 $\geq$8 断言、相同的 2400s 超时、相同的骨干
%（\textbf{doubao-seed-2.0-pro-cloud}；冻结 696-run 快照使用的是 DeepSeek-V4-Pro，现已下线）。
% M1--M8、N1--N2 的结果为 \TBD（待评测）；M0 在最小矩阵中的比较锚点为
% 同 top-120-of-Lite 子集上实测的 M0 Pass@1（\MzeroAnchor{}，12/119 已判定），而非 696-run
% Full-set \MzeroPass{}（后者仍为描述性基线快照）。

\section{Experiments}
\label{sec:experiments}
% 【中文】\section{实验}

{\hbadness=10000\sloppy \textbf{Experimental status.} We formulate 3 core hypotheses (H1--H3) and 7 sub-hypotheses (details in \S\ref{sec:hypotheses}). M0 is \emph{evaluated} on the full \Ncases{}-task set over 3 independent runs (\MzeroPass{}$\pm$\MzeroStd{} Pass@1; DeepSeek-V4-Pro, Cline v3.0.51, \texttt{thinking=low}, 2400s, 16 workers; r1=105/482, r2=96/486, r3=103/490; 3/3 clean, 0 FUSE/quota/burnout). M1--M8 and N1--N2 are to be evaluated. Negative results are reported as publishable boundary conditions.\par}
% 【中文】
% \textbf{实验状态。} 本文提出 3 条核心假设（H1--H3）与 7 条子假设（H1a--c、H2a--c、
% H3）。M0 基线已在完整 \Ncases{} 任务集上\emph{实测}（3 轮独立运行，
% \MzeroPass{}$\pm$\MzeroStd{} Pass@1；DeepSeek-V4-Pro、Cline v3.0.51、\texttt{thinking=low}、
% 2400s 超时、16 workers；r1=105/482、r2=96/486、r3=103/490；3/3 clean、0 FUSE/额度/熔断）。
% 所有其他臂（M1--M8、N1--N2）待评测。阴性结果将作为可发表的边界条件报告。

\textbf{M0 baseline \& backbone comparison.} Table~\ref{tab:m0-baseline} reports the M0 3-run snapshot. Table~\ref{tab:model-comparison} compares 4 backbones on the same \Ncases{}-task set (3 runs each, identical harness). The strongest backbone (DeepSeek-V4-Pro) reaches \MzeroPass{}$\pm$\MzeroStd{}, roughly 2.0$\times$ the weakest viable backbone (doubao-seed-2.0-pro-cloud, \DoubaoPro{}$\pm$\DoubaoProStd{}); a paired McNemar test on majority-vote per-task outcomes confirms the gap is highly significant ($b=69$, $c=2$, $p\approx 0.0000$, $\Delta{=}+10.67$pp).
% 【中文】
% \textbf{M0 基线与骨干对比。} 表~\ref{tab:m0-baseline} 报告 M0 的 3 轮快照。
% 表~\ref{tab:model-comparison} 在同一 \Ncases{} 任务集上对比 4 个骨干（各 3 轮、同一 harness）。
% 最强骨干（DeepSeek-V4-Pro）达到 \MzeroPass{}$\pm$\MzeroStd{}，约为最弱可用骨干
%（doubao-seed-2.0-pro-cloud，\DoubaoPro{}$\pm$\DoubaoProStd{}）的 2.0$\times$；
% 逐题 majority-vote 的 McNemar 配对检验确认差距高度显著
%（$b=69$、$c=2$、$p\approx 0.0000$、$\Delta{=}+10.67$pp）。

\begin{table}[ht]
\centering
\caption{M0 baseline (3-run snapshot on the \Ncases{}-task set, DeepSeek-V4-Pro, 2026-11-24). Pass\% = PASS/(PASS+FAIL). Full per-domain table in Appendix~\ref{sec:additional-results}.}
\label{tab:m0-baseline}
\scriptsize
\setlength{\tabcolsep}{4pt}
\begin{tabular}{@{}lcc@{}}
\toprule
\textbf{Metric} & \textbf{Count} & \textbf{Rate} \\
\midrule
PASS (3 runs total) & 304 & --- \\
FAIL (3 runs total) & 1154 & --- \\
ERROR (3 runs total) & 42 & --- \\
\midrule
Pass@1 (r1) & 105/482 & 21.78\% \\
Pass@1 (r2) & 96/486 & 19.75\% \\
Pass@1 (r3) & 103/490 & 21.02\% \\
Pass@1 (mean$\pm$std) & --- & \textbf{\MzeroPass{}$\pm$\MzeroStd{}} \\
\bottomrule
\end{tabular}
\end{table}

\begin{table}[ht]
\centering
\caption{Backbone comparison on the \Ncases{}-task set (3 runs each, same harness: Cline v3.0.51, \texttt{thinking=low}, 2400s, 16 workers). McNemar test on majority-vote per-task outcomes (DeepSeek-V4-Pro vs doubao-pro): $b=69$, $c=2$, $p\approx 0.0000$, $\Delta{=}+10.67$pp.}
\label{tab:model-comparison}
\scriptsize
\setlength{\tabcolsep}{4pt}
\begin{tabular}{@{}lccc@{}}
\toprule
\textbf{Backbone} & \textbf{Pass@1 (mean$\pm$std)} & \textbf{decided/run} & \textbf{clean} \\
\midrule
DeepSeek-V4-Pro & \textbf{\DSvPro{}$\pm$\DSvProStd{}} & 482/486/490 & 3/3 \\
doubao-seed-2.0-pro & \DoubaoPro{}$\pm$\DoubaoProStd{} & 483/483/484 & 3/3 \\
doubao-seed-2.0-lite & \DoubaoLite{}$\pm$\DoubaoLiteStd{} & 482/484/484 & 3/3 \\
doubao-seed-2.0-mini & \DoubaoMini{}$\pm$\DoubaoMiniStd{} & 483/483/484 & 3/3 \\
\bottomrule
\end{tabular}
\end{table}

\textbf{H1/H2/H3 criteria-check framework.} Each sub-hypothesis's pass/fail is determined by the pre-specified falsification threshold (\S\ref{sec:hypotheses}). Table~\ref{tab:criteria-check} gives the evaluation framework; results are \TBD{} (arms M1--M8, N1--N2 not yet evaluated).
% 【中文】
% \textbf{H1/H2/H3 判据检验框架。} 每条子假设的通过/失败由事前规定的证伪阈值（\S\ref{sec:hypotheses}）
% 决定。表~\ref{tab:criteria-check} 给出评测框架，结果为 \TBD{}（臂 M1--M8、N1--N2 尚未评测）。
% 【中文】表 3：判据检验（结果 \TBD{}，臂未评测）。

\begin{table}[ht]
\centering
\caption{Criteria check (results \TBD{}; arms M1--M8, N1--N2 pending evaluation).}
\label{tab:criteria-check}
\scriptsize
\begin{tabularx}{\columnwidth}{@{}llYcc@{}}
\toprule
\textbf{Hyp} & \textbf{Arm} & \textbf{Test} & \textbf{Est.} & \textbf{Verdict} \\
\midrule
H1a & M3 vs M0 & $\Delta$Pass@1 $\geq$3pp & \TBD & \TBD \\
H1a & M3 vs N1 & exclude prompt-lengthening & \TBD & \TBD \\
H1b & M3+S1 stratified & monotonic $\Delta$ & \TBD & \TBD \\
H1c & 2$\times$2 factorial & interaction $\geq$2pp & \TBD & \TBD \\
H2a & M5 vs M2 & surprise-gate $>$ uniform & \TBD & \TBD \\
H2b & M7 vs M2 & compress $>$ full & \TBD & \TBD \\
H2c & M6 vs M3 & state-aware $>$ agnostic & \TBD & \TBD \\
H3 & single vs spawn & single $\geq$ spawn & \TBD & \TBD \\
\bottomrule
\end{tabularx}
\end{table}

\textbf{Arm-table convergence \& power notes.} The main 10-arm table drops M1/M8 to appendix; M2 is the H2a/H2b control (H1a weak-control criteria withdrawn, \S\ref{sec:arms}). M6's fail-streak test is exploratory at the minimal matrix (H2c deferred to ideal matrix). The minimal-matrix M0 anchor is the 3-run M0 on the full \Ncases{}-task set (\MzeroPass{}$\pm$\MzeroStd{}), not a subset.
% 【中文】
% \textbf{臂表收敛与功效注。} 主 10 臂表将 M1（静态 SOP）与 M8（负效价）移至附录，M2 作为 H2a/H2b 对照
% 运行（非主臂；其 H1a 弱对照判据已撤回，§\ref{sec:arms}）。M6 的失败连续检验在最小矩阵下为探索性
%（H2c 判据延至理想矩阵）。最小矩阵的 M0 比较锚点为完整 \Ncases{} 任务集上的 3 轮 M0
%（\MzeroPass{}$\pm$\MzeroStd{}），而非子集。

\textbf{M0 measurement \& reused assets.} M0 is measured directly on the arm backbone DeepSeek-V4-Pro over the full \Ncases{}-task set: \MzeroPass{}$\pm$\MzeroStd{} Pass@1 (3 runs, r1=105/482, r2=96/486, r3=103/490, 2026-11-24). For H1b stratification, we re-use the trace-length distribution from 47 pilot runs (median 26.5 calls, IQR 13.8--38.0) to define short ($<$14)/mid (14--38)/long ($>$38) strata. M0 is evaluated; M1--M8 and N1--N2 are pending.
% 【中文】
% \textbf{M0 实测与复用资产。} M0 直接在臂骨干 DeepSeek-V4-Pro 上、于完整 \Ncases{} 任务集实测：
% \MzeroPass{}$\pm$\MzeroStd{} Pass@1（3 轮，r1=105/482、r2=96/486、r3=103/490，2026-11-24）。
% 对需要分层的臂（H1b），我们复用 47 条插桩试点 run 的轨迹长度分布
%（中位 26.5 调用，IQR 13.8--38.0）来定义短（$<$14）/中（14--38）/长（$>$38）分层。M0 已评测；
% M1--M8、N1--N2 待评测。

\section{Analysis: Why Memory Helps or Hurts}
\label{sec:analysis}

\textbf{Failure taxonomy (baseline).} Over the \Neval{}-run M0 snapshot, a systematic per-run triage (machine-checkable evidence, explicit UNKNOWN share, never folding unclassified runs) yields: (i) no-output/aborted 21.0\% (146/696), (ii) semantic/partial failure 62.8\% (437/696; 132 runtime crashes + 305 assertion failures), (iii) contract/format 0.3\% (2/696), residual UNKNOWN 7.0\% (49/696). The dominance of semantic failures (62.8\%)---deliverables that are syntactically plausible but semantically wrong---is direct evidence that the binding constraint is \emph{knowledge of what to do}, not output formatting. This supports the memory hypothesis: an agent that consolidated ``which assertions fail and why'' across sessions could convert semantic failures into passes.

\textbf{Trace anatomy (baseline).} 47 instrumented pilot runs (2026-11-03--06) show: median 26.5 tool calls (IQR 13.8--38.0, max 97) across 20.5 iterations, median 1{,}047\,s (IQR 552--1{,}495). Of 1{,}371 tool calls, 41.2\% go through the MCP bridge to live EDA tools and 58.8\% are built-in file/search/shell. The effort is substantial but not unbounded---and the 41.2\% MCP fraction means nearly half of all actions require grounding in live tool outputs the model has never seen during training. (Caveat: these 47 traces are archived pilots, not the \Neval{}-run batch; they describe mechanics, not pass rates.)

\textbf{Affective-mechanism analysis (placeholder, post pre-registration).} The three mechanisms (\S\ref{sec:affective}) predict testable signatures: (1) surprise distribution---M5 should concentrate writes on high-surprise events; (2) $r_{\text{cp}}$ vs.\ Pass@1 scatter---M7 should occupy the ``smaller memory, equal-or-better Pass@1'' region; (3) state\_match hit distribution---M6 should up-weight pitfall memories in the fail-streak subset; (4) M4 leak-rate audit---transfer-hit memories should be abstract cognition, not problem-specific answer fragments (leak rate $=0$). These analyses are \TBD{} pending post-pre-registration evaluation.

\textbf{Why memory might hurt.} If memory is uncurated (M2, no reliability gating), stale or contradictory entries can mislead: a lesson from an old tool version may be wrong for the current one. This is why H1a tests M3 (gated) against M2 (ungated), and why the disconfirmation filter and held-out rollback are part of the reliability gating. If H1a falsifies \emph{but} M3 $\approx$ M2, the negative result itself is publishable: it bounds the value of reliability gating in long-chain flow.

\section{Related Work}
\label{sec:related}

\textbf{EDA benchmarks.} VerilogEval and RTLLM evaluate single-function Verilog generation; FlowFixer-EDA \cite{dong2026flowfixer} targets open-source flow repair. \bench{} is the first to evaluate end-to-end orchestration of multi-step \emph{commercial}-tool design flows from real engineer traces, with $\geq$8 executable assertions and anti-leakage hooks. Our prior script-layer work (DATE27 \cite{PLACEHOLDER_DATE27_ZHULONG}; ZhuLong-DAC \cite{liu2026runreadbackrefine}) established the three-layer memory + RSI framework on short-chain tasks; this paper tests whether those gains transfer to long-chain flow.

\textbf{Agent memory and self-improvement.} Reflexion, ExpeL, and SkillOpt accumulate experience within or across episodes; MUSE provides a create-evaluate-refine loop for skill libraries. Generative Agents \cite{park2023} filter memory writes by an importance score---a functional precursor to our surprise-gated writing (H2a), but without a predictive-error formulation or pre-registered falsification. EM-LLM segments events using Bayesian surprise; PsychoAgent retrieves by affect + conflict. Our contribution is to \emph{pre-register} three affective-memory mechanisms with mechanism-level negative controls (equal-token, mechanism-off) on real long-chain EDA flow, rather than reporting post-hoc gains.

\textbf{Affective neuroscience $\times$ agent memory.} Emotional arousal modulates consolidation via amygdala--hippocampus interaction \cite{mcgaugh2004}; high arousal triggers a central/peripheral encoding trade-off (weapon-focus \cite{loftus1987}); retrieval is mood-congruent and state-dependent \cite{bower1981}; prediction error gates memory formation \cite{ittibaldi2009}. These principles are established in biological cognition; whether they \emph{transfer} to LLM-agent symbolic tool orchestration is the empirical question this paper pre-registers. Full citations in Appendix~\ref{sec:affective-appendix}.

\section{Conclusion: Boundary Conditions of RSI and Affective-Memory Mechanisms}
\label{sec:conclusion}
% 【中文】\section{结论：RSI 与情绪记忆机制的边界条件}

We pre-registered 3 core hypotheses (H1--H3) and 7 sub-hypotheses on whether governed long-term memory + RSI and three affective-memory mechanisms yield transferable gains on long-chain EDA flow. M0 (\MzeroPass{}) anchors the difficulty; other arms are evaluated after \PreregDate.
% 【中文】
% 我们预注册了 3 条核心假设（H1--H3）与 7 条子假设，探讨受治理的长期记忆 + RSI 与三个情绪记忆机制
% 是否在长链 EDA 流程上带来可迁移增益。M0（\MzeroPass{}）锚定难度；其他臂在 \PreregDate 之后评测。

\textbf{Honest route.} If H1/H2 are falsified, this paper reports boundary conditions of RSI / affective heuristics in long-chain flow---both publishable as they delimit where bio-inspired mechanisms apply. If supported, it provides the first pre-registered evidence that affective-memory mechanisms bring mechanistic, not merely functional, gains to long-chain EDA orchestration.
% 【中文】
% \textbf{诚实路线。} 若 H1/H2 被证伪，本文报告长链流程中 RSI / 情绪启发式的边界条件——两者均可发表，
% 因为它们划定了仿生机制的适用范围。若获支持，本文提供首个预注册证据：情绪记忆机制为长链 EDA 编排
% 带来机制性而非纯功能性的增益。

\textbf{Future work.} Multi-modal grounding; longer chains ($>$50 calls); finer valence dimensions; and whether H3 reverses at extreme chain lengths.
% 【中文】
% \textbf{未来工作。} 多模态接地；更长链（$>$50 次调用）；更细的效价维度；以及 H3 是否在极长链长度下反转。


% ===== 参考文献 =====
\clearpage
\bibliographystyle{icml2026}
\bibliography{references}

% ===== Ethics + Reproducibility（ICML 惯例：置于参考文献之后，不计入页数限制）=====
% 编号 10/11：v1.7 正文已扩到 10 节（0_abstract … 9_conclusion），
% 旧版 8_/9_ 前缀与 8_related / 9_conclusion 冲突，故改为 10_/11_。
\section{Ethics Statement}
\label{sec:ethics}

\bench{} tasks are derived from real IC-design engineer coding-agent traces in a production environment. All tasks are desensitized: no proprietary IP, customer data, or identifiable engineer information is released. Commercial EDA tool licenses are not redistributed; the released protocol white-paper and NDA review package allow reviewers to verify claims without access to proprietary tools. Pre-registration discipline (\PreregDate) ensures hypotheses and criteria cannot be retroactively adjusted to favor a result. Negative results are reported as publishable boundary conditions, not suppressed.

\section{Reproducibility Statement}
\label{sec:reproducibility}
% 【中文】\section{可复现性声明}

We release a \emph{protocol white-paper} (harness protocol, evaluation criteria, desensitized task specifications) and an \emph{NDA-controlled review package} granting reviewers access to a sandbox with the commercial EDA tools. We do \textbf{not} release executable Lite-set assets or a mock-EDA mode (v1.4 open-source boundary narrowing). The falsifiability of our claims does not depend on proprietary licenses: an independent party can re-derive an equivalent benchmark on open-source EDA tools (Yosys/OpenROAD/Magic) and test the same 3 core + 7 sub-hypotheses using the released protocol. The M0 baseline (\MzeroPass{}) is reproducible from the frozen 696-run snapshot configuration (DeepSeek-V4-Pro, Cline v3.0.51, \texttt{thinking=low}, 2400s timeout). Full benchmark construction, per-domain results, and experimental details are in Appendices~\ref{sec:benchmark-details}--\ref{sec:experimental-details}.
% 【中文】
% 我们发布\emph{协议白皮书}（执行框架协议、评测判据、脱敏任务规格）与\emph{NDA 受控评审包}，授予
% 评审者访问含商用 EDA 工具的沙箱。我们\textbf{不}发布可执行的 Lite 集资产或 mock-EDA 模式
%（v1.4 开源边界收缩）。我们主张的可证伪性不依赖于专有许可：独立方可在开源 EDA 工具
%（Yosys/OpenROAD/Magic）上重新推导出等价基准，并用发布的协议检验同样的 3 条核心 + 7 条子假设。
% M0 基线（\MzeroPass{}）可从冻结的 696-run 快照配置复现（DeepSeek-V4-Pro、Cline v3.0.51、
% \texttt{thinking=low}、2400s 超时）。完整基准构建、各领域结果与实验细节见
% 附录~\ref{sec:benchmark-details}--\ref{sec:experimental-details}。


% ===== 附录（不计入页数限制）=====
\newpage
\appendix
\section{Benchmark Details}
\label{sec:benchmark-details}
% 【中文】本文件已含 v1.7 时期中文对照注释（`% <chinese>` 行内格式，功能等价；2026-11-18 中文对照令核验保留，未改以守判据③纯新增）。

% 附录 A：benchmark 构建细节（参考 SWE-bench Appendix A "Benchmark Details"）
% 附录 A：基准构建细节（参考 SWE-bench Appendix A "Benchmark Details"）

\subsection{Task Example}
\label{sec:task-example}
% \subsection{任务示例}

[TODO: show one complete task: prompt + assets + check.py + test\_rationale. Suggested: a medium-difficulty schematic\_netlist case that demonstrates the full design\_flow contract.]
% [待办：展示一个完整任务：提示词 + 附件 + check.py + test\_rationale。建议选一个中等难度的 schematic\_netlist 案例，以演示完整的 design\_flow 契约。]

\subsection{Scan Pipeline Details}
\label{sec:scan-pipeline}
% \subsection{扫描流程细节}

We surveyed the engineering workstation fleet and scanned active production machines, yielding several thousand sessions with EDA keyword matches (mc $\geq 15$). Sessions are ranked by match-count (mc)---the number of lines hitting any EDA keyword. The mc $\geq 30$ threshold yields the high-quality candidate pool from which the \emph{Full} set (\FullScale{} tasks, continuously growing, count not reported in the main paper) is curated; a \emph{Lite} subset of \LiteSet{} tasks is filtered from Full by domain-coverage balance (all 13 domains preserved) and chain-length distribution (approximating Full's) for general evaluation under bounded token/time budgets. The exact workstation count, user count, and raw session inventory are company-internal and omitted here; the curation methodology is fully described so that another organization with comparable trace access can replicate the pipeline.
% 我们调研了工程工作站集群并扫描了生产环境中的活跃机器，产出数千个含 EDA 关键词命中的会话（mc $\geq 15$）。会话按匹配计数（mc）——命中任意 EDA 关键词的行数——排序。mc $\geq 30$ 的阈值产出高质量候选池，\emph{Full} 集（\FullScale{} 个任务、持续增长、正文不报告具体数字）即从中整理；\emph{Lite} 子集 \LiteSet{} 个任务则从 Full 按领域覆盖平衡（保留全部 13 域）与链长分布（近似 Full）筛选，用于在有限 token/时间预算下的一般评测。确切的工作站数、用户数与原始会话清单属公司内部信息，此处略去；整理方法论已完整描述，使另一家具有可比 trace 访问权限的组织可以复现该流程。

\textbf{Two-stage prefilter.} A rule-based stage rejects sessions with $< 3$ EDA tool invocations or $> 50\%$ permission-denied lines. A 0.8B LLM classifier (Qwen3.5-0.8B-Q8\_0, CPU) performs a binary is\_design\_flow classification. End-to-end cost is $\sim$1.1s/call at 73 tok/s. On the full candidate pool, the prefilter rejects 0\%---confirming that mc $\geq 30$ sessions are almost universally genuine EDA work.
% \textbf{两级预筛器。} 规则级拒绝 EDA 工具调用 $< 3$ 次或权限被拒行 $> 50\%$ 的会话。一个 0.8B LLM 分类器（Qwen3.5-0.8B-Q8\_0，CPU）执行二分类 is\_design\_flow 判定。端到端开销约 1.1 秒/次，73 tok/s。在完整的候选池上，预筛器拒绝率为 0\%——印证了 mc $\geq 30$ 的会话几乎全为真实 EDA 工作。

\subsection{Domain Keyword Filters}
\label{sec:keyword-filters}
% \subsection{领域关键词过滤器}

Table~\ref{tab:keywords} lists the EDA keywords used for session scanning and area-group assignment. The digital-front group uses an expanded keyword set (synthesis, synopsys, tcl, verilog, systemverilog, rtl, sdc, vivado, yosys, openroad) and a relaxed mc threshold ($\geq 15$) to capture digital-front tasks that use cline less intensively.
% 表~\ref{tab:keywords} 列出用于会话扫描和领域组分配的 EDA 关键词。数字前端组使用扩展关键词集（synthesis、synopsys、tcl、verilog、systemverilog、rtl、sdc、vivado、yosys、openroad）并放宽 mc 阈值（$\geq 15$），以捕获使用 cline 不那么密集的数字前端任务。

\begin{table}[ht]
\centering
\small
\caption{EDA keyword filters by area group.}
% \caption{各领域组的 EDA 关键词过滤器。}
\label{tab:keywords}
\begin{tabular}{@{}llp{0.38\linewidth}@{}}
\toprule
\textbf{Area group} & \textbf{mc threshold} & \textbf{Keywords} \\
\midrule
Digital front & $\geq 15$ & verdi, vcs, dc\_shell, innovus, primetime, synthesis, tcl, verilog, rtl \\
SPICE / char.\ & $\geq 30$ & hspice, finesim, siliconsmart, spectre, ngspice, simulator, spice \\
Physical / layout & $\geq 30$ & calibre, virtuoso, aether, ipcustom, layout, lvs, drc \\
Schematic / analog & $\geq 30$ & waveview, cosmos, skill, netlist, spiros, spectrum, analog, schematic \\
\bottomrule
\end{tabular}
\end{table}

\subsection{Full Domain Distribution}
\label{sec:full-distribution}
% \subsection{完整领域分布}

The benchmark is built in four construction shards. As of 2026-11-21 the living
build holds 2043 cases (A: 443, B: 483, C: 564, D: 554; 36 retired, all in shard A, 2007 active) drawn from the
trace pool described in Section~\ref{sec:scan-pipeline}; these per-shard counts
grow as new traces are curated and reflect the snapshot at submission time (the
main paper reports only the qualitative Full-set scale, not this count). The
Lite subset of \LiteSet{} tasks (used for arm-level evaluation,
Section~\ref{sec:experiments}) is filtered from this living build by
domain-coverage balance and chain-length distribution. The evaluation in
Section~\ref{sec:experiments} uses the frozen \Neval{}-run snapshot of
2026-11-09 (Table~\ref{tab:full-domain-results}). Per-shard counts are maintained
in the benchmark repository; a full user/host distribution table is deferred to
the camera-ready version because the host inventory is company-internal.
% 基准由四个构建分片组成。截至 2026-11-21，活集含 2043 例（A: 443、B: 483、C: 564、
% D: 554；其中 36 例 retired 全属 shard A，active 2007），全部取自第~\ref{sec:scan-pipeline}~节描述的轨迹池；各分片计数随新 trace 整理而增长，
% 此处反映投稿时快照（正文仅报告 Full 集定性规模，不报告此计数）。用于臂级评测的 Lite 子集
%（\LiteSet{} 任务，第~\ref{sec:experiments}~节）即从该活集按领域覆盖平衡与链长分布筛选。
% 第~\ref{sec:experiments}~节的评测使用 2026-11-09 冻结的 \Neval{}-run 快照
%（表~\ref{tab:full-domain-results}）。各分片计数在基准仓库中维护；完整的用户/主机分布表推迟到
% camera-ready，因为主机清单属公司内部信息。
% T12 resolved 2026-11-13：正文中所有工作站/用户/会话级 infra 数字已模糊处理，
% 仅保留方法论描述；Full/Lite 集叙事替代具体扫描统计。

\subsection{Task Schema Specification}
\label{sec:task-schema-spec}
% \subsection{任务模式规范}

Each task is a JSONL record with the following fields:
% 每个任务是一条 JSONL 记录，包含以下字段：
\begin{itemize}
    \item \texttt{id}: unique task identifier (area-group-prefixed, e.g., \texttt{A-001}).
    % \texttt{id}：唯一任务标识（领域组前缀，如 \texttt{A-001}）。
    \item \texttt{mode}: always \texttt{design\_flow}.
    % \texttt{mode}：恒为 \texttt{design\_flow}。
    \item \texttt{prompt}: natural-language task description.
    % \texttt{prompt}：自然语言任务描述。
    \item \texttt{assets}: list of input file paths (copied into the workdir).
    % \texttt{assets}：输入文件路径列表（拷入工作目录）。
    \item \texttt{entry\_point}: \texttt{outputs/submit.zip}.
    % \texttt{entry\_point}：\texttt{outputs/submit.zip}。
    \item \texttt{test}: \texttt{check(deliverable, workdir)} function specification.
    % \texttt{test}：\texttt{check(deliverable, workdir)} 函数规格。
    \item \texttt{test\_rationale}: human-readable assertion list in ``k assertions: (1), (2), ...'' format ($k \geq 8$).
    % \texttt{test\_rationale}：可读断言列表，格式为"k 条断言：(1)、(2)、……"（$k \geq 8$）。
    \item \texttt{trace\_source}: \{host, user, session\_file, session\_msgs, real\_task\_summary\}.
    % \texttt{trace\_source}：\{host, user, session\_file, session\_msgs, real\_task\_summary\}。
    \item \texttt{domain}: one of 13 EDA domains (Section~\ref{sec:full-distribution}).
    % \texttt{domain}：13 个 EDA 领域之一（第~\ref{sec:full-distribution}~节）。
\end{itemize}

\section{Additional Results}
\label{sec:additional-results}
% 【中文】本文件已含 v1.7 时期中文对照注释（`% <chinese>` 行内格式，功能等价；2026-11-18 中文对照令核验保留，未改以守判据③纯新增）。

% 附录 B：扩展实验结果（参考 SWE-bench Appendix C "Additional Results"）
% 附录 B：扩展实验结果（参考 SWE-bench Appendix C "Additional Results"）

\subsection{Full Per-Domain Results}
\label{sec:full-per-domain}
% \subsection{完整的各领域结果}

{\hbadness=10000\sloppy Table~\ref{tab:full-domain-results} reports the per-domain breakdown from the 3-run M0 evaluation (3 runs across all 13 domains, DeepSeek-V4-Pro, Cline v3.0.51, \texttt{thinking=low}, 16 workers, 2400s timeout). The overall Pass@1 is \MzeroPass{}$\pm$\MzeroStd{} (3-run mean$\pm$std; r1=105/482, r2=96/486, r3=103/490 decided; 304 PASS / 1154 FAIL / 42 ERROR-or-ABSENT across 1500 runs). Per-domain breakdown for the 3-run evaluation is pending and will be reported in the camera-ready version; the table below is retained as a placeholder for the per-domain matrix.\par}
% 表~\ref{tab:full-domain-results} 报告 3 轮 M0 评测（3 轮、覆盖全部 13 个领域，
% DeepSeek-V4-Pro、Cline v3.0.51、\texttt{thinking=low}、16 workers、2400s 超时）的各领域分解。
% 总体 Pass@1 为 \MzeroPass{}$\pm$\MzeroStd{}（3-run mean$\pm$std；r1=105/482、r2=96/486、
% r3=103/490 已判定；1500 run 中 304 PASS / 1154 FAIL / 42 ERROR-or-ABSENT）。
% 3-run 评测的逐域分解待补，将在 camera-ready 版本给出；下表保留为逐域矩阵的占位。

\begin{table}[ht]
\centering
\small
\caption{Pass@1 (\%) by domain on the 3-run M0 evaluation (\Ncases{} cases $\times$ 3 runs, DeepSeek-V4-Pro, Cline v3.0.51, \texttt{thinking=low}, 16 workers, 2400s timeout). Per-domain breakdown is pending (TODO); the Total row reports the 3-run aggregate (304 PASS / 1154 FAIL / 42 ERROR-or-ABSENT / 1500 runs; Pass\% = 304/1458 = \MzeroPass{}). Pass\% = PASS/(PASS+FAIL); ERROR and ABSENT are excluded from the denominator.}
% \caption{3 轮 M0 评测（\Ncases{} 例 $\times$ 3 轮，DeepSeek-V4-Pro、Cline v3.0.51、\texttt{thinking=low}、
% 16 workers、超时 2400s）各领域 Pass@1 (\%)。逐域分解待补（TODO）；Total 行报告 3 轮合计
%（304 PASS / 1154 FAIL / 42 ERROR-or-ABSENT / 1500 run；Pass\% = 304/1458 = \MzeroPass{}）。
% Pass\% = PASS/(PASS+FAIL)；ERROR 与 ABSENT 不计入分母。}
\label{tab:full-domain-results}
\scriptsize
\setlength{\tabcolsep}{4pt}
\begin{tabular}{@{}lcccccc@{}}
\toprule
\textbf{Domain} & \textbf{PASS} & \textbf{FAIL} & \textbf{ERROR} & \textbf{ABSENT} & \textbf{Total} & \textbf{Pass\%} \\
\midrule
schematic\_netlist & TODO & TODO & TODO & TODO & TODO & TODO \\
spice\_sim         & TODO & TODO & TODO & TODO & TODO & TODO \\
physical\_verify   & TODO & TODO & TODO & TODO & TODO & TODO \\
digital            & TODO & TODO & TODO & TODO & TODO & TODO \\
analog\_context    & TODO & TODO & TODO & TODO & TODO & TODO \\
waveform           & TODO & TODO & TODO & TODO & TODO & TODO \\
layout\_virtuoso   & TODO & TODO & TODO & TODO & TODO & TODO \\
layout\_aether     & TODO & TODO & TODO & TODO & TODO & TODO \\
verification       & TODO & TODO & TODO & TODO & TODO & TODO \\
eda\_other         & TODO & TODO & TODO & TODO & TODO & TODO \\
pg\_layout         & TODO & TODO & TODO & TODO & TODO & TODO \\
characterization   & TODO & TODO & TODO & TODO & TODO & TODO \\
pnr\_signoff       & TODO & TODO & TODO & TODO & TODO & TODO \\
\midrule
\textbf{Total}     & \textbf{304} & \textbf{1154} & \multicolumn{2}{c}{\textbf{42}} & \textbf{1500} & \textbf{\MzeroPass{}} \\
\bottomrule
\end{tabular}
\end{table}

\subsection{Per-Model Results}
\label{sec:per-model}
% \subsection{各模型结果}

This submission reports a \emph{four-backbone} comparison on the \Ncases{}-task set (3 runs each, Pass@1, same harness: Cline v3.0.51, \texttt{thinking=low}, 2400s timeout, 16 workers). DeepSeek-V4-Pro reaches \DSvPro{}$\pm$\DSvProStd{} (r1=21.78\%, r2=19.75\%, r3=21.02\%; 3/3 runs clean), doubao-seed-2.0-pro-cloud \DoubaoPro{}$\pm$\DoubaoProStd{}, doubao-seed-2.0-lite-cloud \DoubaoLite{}$\pm$\DoubaoLiteStd{}, doubao-seed-2.0-mini-cloud \DoubaoMini{}$\pm$\DoubaoMiniStd{}. A McNemar paired test on majority-vote outcomes (DeepSeek-V4-Pro vs doubao-seed-2.0-pro-cloud) gives $b{=}69$, $c{=}2$, $p\approx 0.0000$ ($\Delta{=}{+}10.67$pp, highly significant), confirming that the M0 gap is not noise. Additional backbones (GLM-5.2, Kimi-K2.6, DeepSeek-V4-Flash) are configured in the same harness but their sweeps are still running; we will report them in the camera-ready version as the multi-backbone test of RQ3 (Section~\ref{sec:per-model}).
% 本投稿给出 \Ncases{}-任务集上的\emph{四骨干}对比（各 3 轮、Pass@1、同一框架：Cline v3.0.51、
% \texttt{thinking=low}、2400s 超时、16 workers）。DeepSeek-V4-Pro 达 \DSvPro{}$\pm$\DSvProStd{}
%（r1=21.78\%、r2=19.75\%、r3=21.02\%；3/3 runs clean），doubao-seed-2.0-pro-cloud \DoubaoPro{}$\pm$\DoubaoProStd{}，
% doubao-seed-2.0-lite-cloud \DoubaoLite{}$\pm$\DoubaoLiteStd{}，doubao-seed-2.0-mini-cloud
% \DoubaoMini{}$\pm$\DoubaoMiniStd{}。对 majority-vote 结果做 McNemar 配对检验
%（DeepSeek-V4-Pro vs doubao-seed-2.0-pro-cloud）得 $b{=}69$、$c{=}2$、$p\approx 0.0000$
%（$\Delta{=}{+}10.67$pp，高度显著），印证 M0 差距非噪声。其余骨干（GLM-5.2、Kimi-K2.6、
% DeepSeek-V4-Flash）已接入同一框架但扫描仍在跑；camera-ready 版本将给出该多骨干检验作为 RQ3。

\subsection{Failure Taxonomy by Domain}
\label{sec:failure-by-domain}
% \subsection{各领域失败分类法}

{\hbadness=10000\sloppy
The batch runner records one terminal status per run, so the per-domain failure
\emph{categories} cannot be derived from that terminal status alone. We therefore
report (i) the run-level outcome anchors per domain (Table~\ref{tab:full-domain-results})
and (ii) a triage procedure that recovers the categories from each run's judge log and
agent session (\texttt{eval\_triage/05\_batch\_triage.py}, available to qualified reviewers
via the NDA-gated review package rather than publicly released, with an explicit UNKNOWN
share so that unclassified runs are never silently folded into another bucket). Applying it to
the 3-run \Ncases{}-task batch yields the per-domain category matrix in
Table~\ref{tab:failure-by-domain} (PASS runs are out of scope here; rows sum to
the FAIL+ERROR count per domain). Class abbreviations: TO=ERROR\_TIMEOUT,
INFRA=ERROR\_INFRA, NOSUB=NO\_SUBMIT, FMT=FORMAT\_CONTRACT, RT=ASSERT\_RUNTIME,
SEM=ASSERT\_SEMANTIC, ANTI=ASSERT\_ANTIHARDCODE, UNK=UNKNOWN; F+E=FAIL+ERROR.
The per-domain category counts for the 3-run batch are pending (TODO); a stacked-bar view and a Cohen's $\kappa$ agreement
statistic between the tool and a human labeler on a 100-run sample are deferred
to the camera-ready version.
\par}
% 批次 runner 每个 run 只记录一个终态，因此仅凭该终态无法推导各领域的失败\emph{类别}。
% 故我们给出：(i) 各领域 run 级结果锚点（表~\ref{tab:full-domain-results}）；
% (ii) 一个逐 run 归因流程，它从每个 run 的判分日志与智能体会话中恢复类别
% （eval_triage/05_batch_triage.py，经 NDA 受控评审包向合格评审者提供，不公开发布；
% 并显式报告 UNKNOWN 比例，避免未分类 run 被悄悄并入别的桶）。
% 将其应用于 3 轮 \Ncases{}-任务批次得到表~\ref{tab:failure-by-domain} 的逐域类别矩阵（PASS run 不在此表
% 范围；各行加和为该域 FAIL+ERROR 数）。列缩写：TO=ERROR_TIMEOUT, INFRA=ERROR_INFRA,
% NOSUB=NO_SUBMIT, FMT=FORMAT_CONTRACT, RT=ASSERT_RUNTIME, SEM=ASSERT_SEMANTIC,
% ANTI=ASSERT_ANTIHARDCODE, UNK=UNKNOWN; F+E=FAIL+ERROR。3 轮批次的逐域类别计数待补（TODO）；
% 堆叠柱状图与工具对人工标注者在 100-run 样本上的 Cohen's kappa 一致性推迟到 camera-ready。

\begin{table}[ht]
\centering
\caption{Per-domain failure-category matrix for the 3-run \Ncases{}-task batch, recovered by \texttt{eval\_triage/05\_batch\_triage.py} from per-run judge logs. PASS runs are excluded; each row sums to that domain's FAIL+ERROR count. UNKNOWN is reported as its own column, not folded into another class. Per-domain category counts are pending (TODO); overall Pass\% is \MzeroPass{} (304/1458 decided across 3 runs).}
\label{tab:failure-by-domain}
\resizebox{\columnwidth}{!}{%
\begin{tabular}{lrrrrrrrrr}
\toprule
Domain & TO & INFRA & NOSUB & FMT & RT & SEM & ANTI & UNK & F+E \\
\midrule
schematic\_netlist & TODO & TODO & TODO & TODO & TODO & TODO & TODO & TODO & TODO \\
spice\_sim         & TODO & TODO & TODO & TODO & TODO & TODO & TODO & TODO & TODO \\
physical\_verify   & TODO & TODO & TODO & TODO & TODO & TODO & TODO & TODO & TODO \\
digital            & TODO & TODO & TODO & TODO & TODO & TODO & TODO & TODO & TODO \\
analog\_context    & TODO & TODO & TODO & TODO & TODO & TODO & TODO & TODO & TODO \\
waveform           & TODO & TODO & TODO & TODO & TODO & TODO & TODO & TODO & TODO \\
layout\_virtuoso   & TODO & TODO & TODO & TODO & TODO & TODO & TODO & TODO & TODO \\
layout\_aether     & TODO & TODO & TODO & TODO & TODO & TODO & TODO & TODO & TODO \\
verification       & TODO & TODO & TODO & TODO & TODO & TODO & TODO & TODO & TODO \\
eda\_other         & TODO & TODO & TODO & TODO & TODO & TODO & TODO & TODO & TODO \\
pg\_layout         & TODO & TODO & TODO & TODO & TODO & TODO & TODO & TODO & TODO \\
characterization   & TODO & TODO & TODO & TODO & TODO & TODO & TODO & TODO & TODO \\
pnr\_signoff       & TODO & TODO & TODO & TODO & TODO & TODO & TODO & TODO & TODO \\
\midrule
\textbf{Total} & \textbf{TODO} & \textbf{TODO} & \textbf{TODO} & \textbf{TODO} & \textbf{TODO} & \textbf{TODO} & \textbf{TODO} & \textbf{TODO} & \textbf{TODO} \\
\bottomrule
\end{tabular}}
\end{table}

\subsection{Token Usage and Cost Analysis}
\label{sec:token-cost}
% \subsection{Token 用量与成本分析}

We can report \emph{behavioural} cost from the archived pilot sessions. For the 30
runs that carry an authoritative \texttt{run\_result.aggregateUsage} record, the
median run consumes 0.71M input and 55.6k output tokens (mean 0.86M/52.5k; totals
25.7M/1.58M), on top of a median of 26.5 tool calls (IQR 13.8--38.0) over 20.5
iterations and 1{,}047\,s of wall-clock time (IQR 552--1{,}495). The input:output
ratio of ${\sim}13{:}1$ is the signature of a long-horizon agent loop that re-sends
its context each iteration, and it is the main cost driver of flow orchestration.
A per-task token/cost table for the 3-run \Ncases{}-task batch is pending because the
batch summary records only terminal status and wall time; the per-run session logs
are retained and will be aggregated with the same parser
(\texttt{eval\_triage/trace\_anatomy.py}) in the camera-ready version.
% 我们能报告的是\emph{行为学}成本。对 30 份带有权威 run_result.aggregateUsage 记录的 run：
% 中位每 run 消耗 0.71M input / 55.6k output token（均值 0.86M / 52.5k；合计 25.7M / 1.58M），
% 同时中位发出 26.5 次工具调用（IQR 13.8--38.0）、跨 20.5 轮迭代、耗 1,047 秒墙钟（IQR 552--1,495）。
% 约 13:1 的 input:output 比是长时程智能体回路的特征——每轮迭代都要重发上下文，这也是流程编排的主要成本来源。
% 3 轮 \Ncases{}-任务批次的逐任务 token/成本表待补：批次摘要只记录终态与墙钟，逐 run 会话日志已留存，
% 将在 camera-ready 用同一解析器（eval_triage/trace_anatomy.py）聚合给出。

\subsection{Context-Compaction Ablation}
\label{sec:compaction-ablation}
% \subsection{上下文压缩消融}

Long-horizon agent loops accumulate context until the coding agent's compaction
mechanism must summarize earlier turns to fit the context window. Cline offers
two compaction strategies: \emph{basic} (a structured summarization that
preserves tool-call I/O verbatim) and \emph{agentic} (an LLM-generated free-form
summary). Because compaction is triggered deep into a run and its effect is
entangled with task difficulty, we isolate it with an ABBA-designed ablation on
a fixed 40-task subset of the \Ncases{}-task benchmark, holding everything else constant
(DeepSeek-V4-Pro, Cline v3.0.51, \texttt{thinking=medium}, 2400s timeout, single
trace per task). The ABBA order (basic, agentic, agentic, basic, basic, agentic)
balances position effects across six rounds.

\begin{table}[ht]
\centering
\small
\caption{Context-compaction ablation (ABBA, 6 rounds, 40-task subset). Pass\% is
PASS/40. Basic compaction consistently outperforms agentic by ${\sim}10$pp while
keeping wall-clock time lower, suggesting that preserving tool-call I/O verbatim
is more valuable than a lossy semantic summary for EDA flow orchestration.}
\label{tab:compaction-ablation}
\footnotesize
\setlength{\tabcolsep}{4pt}
\begin{tabular}{@{}llcccc@{}}
\toprule
\textbf{Round} & \textbf{Variant} & \textbf{Pass\%} & \textbf{PASS/40} & \textbf{Wall (h)} & \textbf{Tok/task (M)} \\
\midrule
r1 & basic   & 67.5 & 27 & 1.54 & 2.59 \\
r1 & agentic & 57.5 & 23 & 2.39 & 1.90 \\
r2 & agentic & 60.0 & 24 & 1.88 & 2.03 \\
r2 & basic   & 70.0 & 28 & 1.83 & 2.06 \\
r3 & basic   & 70.0 & 28 & 1.80 & 2.03 \\
r3 & agentic & 60.0 & 24 & 1.93 & 1.82 \\
\midrule
\multicolumn{2}{l}{basic mean}   & 69.2 & 27.7 & 1.72 & 2.23 \\
\multicolumn{2}{l}{agentic mean} & 59.2 & 23.7 & 2.07 & 1.92 \\
\bottomrule
\end{tabular}
\end{table}

Across the three basic rounds the pass rate is 67.5/70.0/70.0 (mean 69.2\%) and
across the three agentic rounds 57.5/60.0/60.0 (mean 59.2\%), a stable
${\sim}10$pp gap in favour of basic compaction. The ABBA design rules out a
pure ordering artifact, and the gap is consistent in direction across all three
round pairs. Interestingly, basic compaction is also \emph{faster} (1.72\,h vs
2.07\,h mean batch wall time) despite consuming slightly more tokens (2.23\,M vs
1.92\,M per task): the agentic summarizer spends extra iterations re-deriving
information that the basic summarizer keeps verbatim. We interpret this as
evidence that in EDA flow orchestration the high-value context is tool-call I/O
(exact commands, exact tool responses, exact file paths), which a free-form
semantic summary tends to lose. The ablation uses \texttt{thinking=medium}
rather than the \texttt{thinking=low} of the main evaluation, so the absolute
pass rates are not directly comparable to Table~\ref{tab:full-domain-results};
the relative comparison between the two compaction strategies is the point.
% 长时程智能体回路不断累积上下文，直到编程智能体的压缩机制必须摘要早期轮次以适配上下文窗口。
% Cline 提供两种压缩策略：\emph{basic}（保留工具调用 I/O 原文的结构化摘要）与 \emph{agentic}
% （LLM 生成的自由式摘要）。由于压缩在 run 深处触发且其效果与任务难度耦合，我们用 ABBA 设计的消融
% 在 \Ncases{}-任务基准的固定 40 任务子集上隔离该因素，其余保持恒定（DeepSeek-V4-Pro、Cline v3.0.51、
% \texttt{thinking=medium}、2400s 超时、每任务单轨迹）。ABBA 顺序（basic, agentic, agentic, basic,
% basic, agentic）在六轮间平衡位置效应。三轮 basic 的通过率为 67.5/70.0/70.0（均值 69.2\%），
% 三轮 agentic 为 57.5/60.0/60.0（均值 59.2\%），basic 稳定优 ${\sim}10$pp。ABBA 设计排除了纯顺序假象，
% 且三轮对的方向一致。有趣的是 basic 同时更快（均值 1.72h vs 2.07h）但略多耗 token（2.23M vs 1.92M/任务）：
% agentic 摘要器需额外迭代重新推导 basic 摘要器保留为原文的信息。我们将此解读为：在 EDA 流程编排中，
% 高价值上下文是工具调用 I/O（精确命令、精确工具响应、精确文件路径），自由式语义摘要倾向于丢失这些。
% 该消融使用 \texttt{thinking=medium} 而非主评测的 \texttt{thinking=low}，故绝对通过率与
% 表~\ref{tab:full-domain-results} 不可直接比较；两种压缩策略的相对比较才是重点。

\section{Experimental Details}
\label{sec:experimental-details}
% 【中文】\section{实验细节}
% 附录 C：实验细节（参考 SWE-bench Appendix D "Additional Experimental Details"）
% 注：本文件原含 GBK→UTF-8 双重解码导致的乱码中文注释，2026-11-18 中文对照令期间修正为正确 UTF-8。

\subsection{Cline Agent Configuration}
\label{sec:cline-config}
% 【中文】\subsection{Cline 智能体配置}

\textbf{Version.} Cline v3.0.51~\cite{cline2026}, open-source.
% 【中文】\textbf{版本。} Cline v3.0.51~\cite{cline2026}，开源。

\textbf{Thinking budget.} \texttt{thinking=low} (default). We chose low over medium/high to control token cost across the \LiteSet{}-task Lite evaluation set; the trade-off is reduced reasoning depth. A separate context-compaction ablation (basic vs.\ agentic summarization, \texttt{thinking=medium}, 40-task subset) is reported in Appendix~\ref{sec:compaction-ablation}; a full thinking-budget sweep is planned for the camera-ready version.
% 【中文】
% \textbf{思考预算。} \texttt{thinking=low}（默认）。我们选 low 而非 medium/high，以在 \LiteSet{}-任务的
% Lite 评测集上控制 token 成本；代价是推理深度降低。一个单独的上下文压缩消融（basic vs.\ agentic 摘要，
% \texttt{thinking=medium}、40 任务子集）见附录~\ref{sec:compaction-ablation}；完整的思考预算扫描计划在
% camera-ready 版本进行。

\textbf{Auto-approve.} All tool calls auto-approved (no human-in-the-loop). The agent may iterate freely within the 2400s timeout.
% 【中文】
% \textbf{自动批准。} 所有工具调用自动批准（无人工介入）。智能体可在 2400 秒超时内自由迭代。

{\hbadness=10000\sloppy \textbf{MCP bridge.} The \texttt{eda\_flow} MCP server (port 18891) exposes 8 tools: \texttt{open\_session}, \texttt{close\_session}, \texttt{declare\_done}, \texttt{run\_commands}, \texttt{run\_code}, \texttt{read\_file}, \texttt{write\_file}, \texttt{list\_files}. The agent must call \texttt{open\_session} first to obtain a workdir, then pass it to all subsequent calls.\par}
% 【中文】
% \textbf{MCP 网桥。} \texttt{eda\_flow} MCP 服务器（端口 18891）暴露 8 个工具：\texttt{open\_session}、
% \texttt{close\_session}、\texttt{declare\_done}、\texttt{run\_commands}、\texttt{run\_code}、
% \texttt{read\_file}、\texttt{write\_file}、\texttt{list\_files}。智能体必须先调用
% \texttt{open\_session} 取得工作目录，再将其传入所有后续调用。

\subsection{API Gateway Configuration}
\label{sec:api-config}
% 【中文】\subsection{API 网关配置}

All models accessed via an OpenAI-compatible API gateway (\texttt{agi-gateway}). The \texttt{providers.json} configuration (with API keys masked):
% 【中文】
% 所有模型通过 OpenAI 兼容 API 网关（\texttt{agi-gateway}）访问。\texttt{providers.json} 配置
%（API 密钥已脱敏）：
{\scriptsize\begin{verbatim}
{
  "baseURL": "http://agi-gateway.internal/cloud/v1",
  "apiKey": "<MASKED>",
  "model": "doubao-seed-2.0-pro-cloud"
}
\end{verbatim}}
% 注：主机名已匿名化为 agi-gateway.internal（双盲投稿，原文含公司内部域名）。
% 注：arm backbone 为 doubao-seed-2.0-pro-cloud（696-run 描述性快照用 DeepSeek-V4-Pro，已下线）。
{\hbadness=10000\sloppy We verify API reachability with a \texttt{/chat/completions} health check before each batch run; a 403 response indicates key expiry.\par}
% 【中文】
% 每批运行前，我们用 \texttt{/chat/completions} 健康检查验证 API 可达性；403 响应表示密钥失效。

\subsection{Anti-Leakage Hook Implementation}
\label{sec:hook-impl}
% 【中文】\subsection{防泄漏钩子实现}

The \texttt{PreToolUse} hook is a shell script invoked before every tool call. It enforces:
% 【中文】
% \texttt{PreToolUse} 钩子是一个 shell 脚本，在每次工具调用前触发。它强制执行：
{\hbadness=10000\sloppy\begin{itemize}
    \item \textbf{Path whitelist}: only paths under the evaluation workdir root are accessible.
    % 【中文】\textbf{路径白名单}：仅评测工作目录根之下的路径可访问。
    \item \textbf{Answer-key blocklist}: any path matching \texttt{check.py}, \texttt{\_check.py}, \texttt{reference/}, \texttt{hidden/}, or \texttt{\_reference\_check.py} is blocked.
    % 【中文】\textbf{答案黑名单}：任何匹配 \texttt{check.py}、\texttt{\_check.py}、\texttt{reference/}、\texttt{hidden/}、\texttt{\_reference\_check.py} 的路径被阻止。
    \item \textbf{Tool whitelist}: only \texttt{read\_file}, \texttt{write\_to\_file}, \texttt{eda\_flow/*} MCP tools, and basic shell commands are allowed; code-search tools (\texttt{grep}, \texttt{find} outside the workdir) are blocked.
    % 【中文】\textbf{工具白名单}：仅允许 \texttt{read\_file}、\texttt{write\_to\_file}、\texttt{eda\_flow/*} MCP 工具及基本 shell 命令；工作目录之外的 \texttt{grep}、\texttt{find} 被阻止。
    \item \textbf{Blocked-call response}: the hook returns the string ``ACCESS RESTRICTED'' to the agent, which is logged but does not terminate the session.
    % 【中文】\textbf{被拦截调用的响应}：钩子向智能体返回字符串"ACCESS RESTRICTED"，该事件被记录但不终止会话。
\end{itemize}}
We count blocked calls per session (\texttt{hook\_blocks} field in \texttt{status.jsonl}) to verify the hook is firing.
% 【中文】
% 我们统计每会话被拦截次数（\texttt{status.jsonl} 的 \texttt{hook\_blocks} 字段），以验证钩子是否触发。

\subsection{Judge Implementation}
\label{sec:judge-impl}
% 【中文】\subsection{判分器实现}

The judge runs \emph{after} the agent exits, in a separate directory:
% 【中文】
% 判分器在智能体退出\emph{之后}运行，位于独立目录：
\begin{enumerate}
    \item Unpack \texttt{submit.zip} into a fresh workdir.
    % 【中文】将 \texttt{submit.zip} 解包到全新工作目录。
    \item Invoke \texttt{check.py} with \texttt{(deliverable\_path, workdir)} arguments.
    % 【中文】以 \texttt{(deliverable\_path, workdir)} 参数调用 \texttt{check.py}。
    \item \texttt{check.py} uses only Python standard library (no EDA license required).
    % 【中文】\texttt{check.py} 仅使用 Python 标准库（判分无需 EDA 许可）。
    \item A task passes iff all $k \geq 8$ assertions hold.
    % 【中文】当且仅当所有 $k \geq 8$ 条断言成立时，任务判为通过。
\end{enumerate}
The \texttt{test\_rationale} field documents each assertion in human-readable form (``k assertions: (1), (2), ...'') for interpretability and audit.
% 【中文】
% \texttt{test\_rationale} 字段以可读形式（"k 条断言：(1)、(2)、……"）记录每条断言，便于可解释性与审计。

\subsection{Open-Source EDA Equivalence Mapping}
\label{sec:repro-details}
% 【中文】\subsection{开源 EDA 等价映射}

To address the reproducibility limitation of commercial EDA tools, an independent party can construct an open-source-EDA equivalent by mapping each domain to an open-source tool:
% 【中文】
% 为应对商用 EDA 工具的可复现性限制，独立方可通过将每个领域映射到对应开源工具来构建开源 EDA 等价基准：
{\hbadness=10000\sloppy\begin{itemize}
    \item \textbf{Digital synthesis}: Synopsys DC $\to$ Yosys + OpenROAD.
    % 【中文】\textbf{数字综合}：Synopsys DC $\to$ Yosys + OpenROAD。
    \item \textbf{SPICE simulation}: HSPICE/Spectre $\to$ ngspice/Xyce.
    % 【中文】\textbf{SPICE 仿真}：HSPICE/Spectre $\to$ ngspice/Xyce。
    \item \textbf{Physical verification}: Calibre $\to$ KLayout + Magic.
    % 【中文】\textbf{物理验证}：Calibre $\to$ KLayout + Magic。
    \item \textbf{Layout}: Virtuoso/Aether $\to$ KLayout + OpenROAD.
    % 【中文】\textbf{版图}：Virtuoso/Aether $\to$ KLayout + OpenROAD。
\end{itemize}}
Such an open-source re-implementation would reuse the task schema, harness, and anti-leakage protocol described in this paper, enabling cross-ecosystem comparison. We do not release such an extension; the mapping is provided so that independent parties can re-derive an equivalent benchmark.
% 【中文】
% 此类开源再实现将复用本论文描述的任务模式、框架与防泄漏协议，从而支持跨生态比对。我们不发布该扩展；
% 提供此映射以便独立方可重新推导出等价基准。

\section{Affective-Memory Mechanism Design Details}
\label{sec:affective-appendix}
% 【中文】\section{情绪记忆机制设计细节}

This appendix gives full derivations of the three affective-memory mechanisms introduced in \S\ref{sec:affective}, the cognitive-neuroscience evidence, and the AI-affective-memory literature positioning. All citations were verified via the research.cimi\_search service (8 calls); BibTeX entries are in \texttt{main.bbl}.
% 【中文】
% 本附录给出 \S\ref{sec:affective} 引入的三个情绪记忆机制的完整推导、认知神经科学证据，以及 AI 情绪
% 记忆文献定位。所有引用经 research.cimi\_search 服务核验（8 次调用）；BibTeX 条目见 \texttt{main.bbl}。

\subsection{Mechanism 1: Surprise/Penalty-Gated Writing}
% 【中文】\subsection{机制 1：惊异/惩罚门控写入}

\textbf{Biological evidence.} Emotional arousal---regardless of valence---enhances memory consolidation via the amygdala--hippocampus interaction: adrenal epinephrine/cortisol $\rightarrow$ amygdala norepinephrine $\rightarrow$ hippocampal LTP \cite{mcgaugh2004}. Richter-Levin \& Akirav (2003) describe ``emotional tagging'' of salient events.
% 【中文】
% \textbf{生物证据。} 情绪觉醒——不分效价——经杏仁核-海马交互增强记忆巩固：肾上腺素/皮质醇 $\rightarrow$
% 杏仁核去甲肾上腺素 $\rightarrow$ 海马 LTP \cite{mcgaugh2004}。Richter-Levin \& Akirav (2003) 描述了
% 显著事件的"情绪标记"。

\textbf{Computational mapping.} The write-gating weight (Eq.~\ref{eq:write-gate}) uses Bayesian surprise $\text{surprise}(e) = \text{KL}(p_{\text{posterior}} \| p_{\text{prior}})$ as the prediction-error signal \cite{ittibaldi2009}, and penalty as the sunk cost of a mid-chain failure (tool calls / tokens consumed before the crash). The agent writes an event $e$ to L1/L2 with priority $\propto w_{\text{write}}(e)$, subject to a total token budget held equal to the M2 uniform-sampling baseline. Calibrating $\alpha,\beta$: M5 sweeps $(\alpha,\beta) \in [0,1]^2$ on the transfer target domain; the Pass@1 iso-contour identifies the operating point. \textbf{Falsification (H2a):} M5 $-$ M2 $<$ 1pp or direction reversal $\Rightarrow$ surprise-gating has no advantage over uniform sampling at equal token budget.
% 【中文】
% \textbf{可计算映射。} 写入门控权重（Eq.~\ref{eq:write-gate}）以贝叶斯惊异
% $\text{surprise}(e) = \text{KL}(p_{\text{posterior}} \| p_{\text{prior}})$ 作为预测误差信号
% \cite{ittibaldi2009}，以惩罚作为链中途失败的沉没成本（崩溃前消耗的工具调用/token）。智能体以
% $\propto w_{\text{write}}(e)$ 的优先级将事件 $e$ 写入 L1/L2，受总 token 预算约束（与 M2 均匀采样
% 基线相等）。$\alpha,\beta$ 校准：M5 在迁移目标域上扫描 $(\alpha,\beta) \in [0,1]^2$；Pass@1 等值线
% 确定工作点。\textbf{证伪（H2a）：} M5 $-$ M2 $<$ 1pp 或方向反转 $\Rightarrow$ 惊异门控在等 token
% 预算下相对均匀采样无优势。

\subsection{Mechanism 2: Weapon-Focus Compression}
% 【中文】\subsection{机制 2：武器焦点压缩}

\textbf{Biological evidence.} Under high arousal, the brain consolidates \emph{central} details (the threat) at the expense of \emph{peripheral} details---the ``weapon-focus effect'' \cite{loftus1987}; Easterbrook (1959) attentional narrowing; Mather \& Sutherland (2011) arousal-biased competition.
% 【中文】
% \textbf{生物证据。} 高觉醒下，大脑以\emph{外围}细节为代价巩固\emph{中心}细节（威胁）——"武器焦点
% 效应"\cite{loftus1987}；Easterbrook (1959) 注意力收窄；Mather \& Sutherland (2011) 觉醒偏向竞争。

\textbf{Computational mapping.} The compression ratio $r_{\text{cp}} = |\text{central tokens}|/|\text{total tokens}|$ (Eq.~\ref{eq:weapon-focus}) is triggered when surprise exceeds the same threshold as M5. \emph{Central} = error stack trace + root-cause causal chain (extracted by a root-cause abstraction step); \emph{peripheral} = intermediate tool-call logs. M7 writes only central content; M2 writes everything. \textbf{Falsification (H2b):} M7 $<$ M2 $-$ 2pp (compression lost critical info) or M7 memory $\geq$ 0.8$\times$M2 (compression ineffective) $\Rightarrow$ weapon-focus compression is harmful or inert in EDA flow.
% 【中文】
% \textbf{可计算映射。} 压缩比 $r_{\text{cp}} = |\text{central tokens}|/|\text{total tokens}|$
%（Eq.~\ref{eq:weapon-focus}）在惊异超过与 M5 相同阈值时触发。\emph{中心} = 错误堆栈轨迹 + 根因因果链
%（由根因抽象步骤提取）；\emph{外围} = 中间工具调用日志。M7 仅写中心内容；M2 写全部。
% \textbf{证伪（H2b）：} M7 $<$ M2 $-$ 2pp（压缩丢失关键信息）或 M7 记忆 $\geq$ 0.8$\times$M2
%（压缩无效）$\Rightarrow$ 武器焦点压缩在 EDA 流程中有害或惰性。

\subsection{Mechanism 3: State-Aware Retrieval}
% 【中文】\subsection{机制 3：状态感知检索}

\textbf{Biological evidence.} Bower (1981) Network Theory of Affect: mood-congruent recall (matching valence) and state-dependent recall (matching encoding/retrieval state) \cite{bower1981}; Eich et al.\ (1994) show stronger effects under self-generated events.
% 【中文】
% \textbf{生物证据。} Bower (1981) 情绪网络理论：情绪一致回忆（匹配效价）与状态依赖回忆（匹配编码/检索
% 状态）\cite{bower1981}；Eich et al.\ (1994) 显示在自我生成事件下效应更强。

\textbf{Computational mapping.} The retrieval score (Eq.~\ref{eq:state-retrieval}) adds a state-match term to semantic similarity: $\text{state} = (\text{recent-}N\text{-step failure rate}, \text{current tool domain}, \text{chain stage})$. An agent in a fail-streak retrieves pitfall-avoidance memories; in smooth progress, it retrieves success templates. M6 uses state-aware retrieval on M5-sourced memory (controlling the write side); M3 uses state-agnostic (sim-only) retrieval. \textbf{Falsification (H2c):} fail-streak subset M6 $-$ M3 $<$ 1pp or reversal, or smooth-progress subset M6 $<$ M3 $-$ 2pp (state-awareness harms in good conditions) $\Rightarrow$ H2c falsified.
% 【中文】
% \textbf{可计算映射。} 检索评分（Eq.~\ref{eq:state-retrieval}）在语义相似度上加入状态匹配项：
% $\text{state} = (\text{最近 }N\text{ 步失败率}, \text{当前工具领域}, \text{链阶段})$。处于失败连击的
% 智能体检索避坑记忆；顺利推进时检索成功模板。M6 对 M5 来源记忆使用状态感知检索（控制写入侧）；M3 使用
% 状态无关（仅 sim）检索。\textbf{证伪（H2c）：} 失败连击子集 M6 $-$ M3 $<$ 1pp 或反转，或顺利子集
% M6 $<$ M3 $-$ 2pp（状态感知在好条件下有害）$\Rightarrow$ H2c 被证伪。

\subsection{AI-Affective-Memory Literature Positioning}
% 【中文】\subsection{AI 情绪记忆文献定位}

\begin{itemize}[noitemsep,topsep=0pt]
    \item \textbf{Generative Agents} \cite{park2023}: importance-score-filtered writes---a functional precursor to H2a, but without a predictive-error formulation or mechanism-level falsification.
    % 【中文】\textbf{Generative Agents} \cite{park2023}：按重要性分数过滤写入——H2a 的功能性前驱，但无预测误差形式化或机制级证伪。
    \item \textbf{EM-LLM} (arXiv:2407.09450): segments events using Bayesian surprise---shares the surprise signal but targets episodic segmentation, not write-gating for tool orchestration.
    % 【中文】\textbf{EM-LLM}（arXiv:2407.09450）：用贝叶斯惊异切分事件——共享惊异信号，但面向情节切分，而非工具编排的写入门控。
    \item \textbf{PsychoAgent} (arXiv:2608.07438): retrieves by affect + conflict---a functional precursor to H2c, but without state-dependent encoding or mechanism-level negative controls.
    % 【中文】\textbf{PsychoAgent}（arXiv:2608.07438）：按情绪 + 冲突检索——H2c 的功能性前驱，但无状态依赖编码或机制级阴性对照。
    \item \textbf{Somatic-marker re-injection} (arXiv:2605.08611): re-injects emotion vectors---orthogonal to our write-gate/compression/retrieval decomposition.
    % 【中文】\textbf{躯体标记再注入}（arXiv:2605.08611）：再注入情绪向量——与我们的写入门控/压缩/检索分解正交。
\end{itemize}

Our contribution: \emph{first falsifiable test} of these three bio-inspired mechanisms on real long-chain EDA flow, with mechanism-level negative controls (equal-token, mechanism-off) rather than post-hoc gain reports.
% 【中文】
% 我们的贡献：在真实长链 EDA 流程上对这些仿生机制的\emph{首个可证伪检验}，配以机制级阴性对照
%（等 token、机制关闭），而非事后增益报告。

\subsection{Placeholder Figures (pending evaluation)}
% 【中文】\subsection{占位图（待评测）}

\begin{itemize}[noitemsep,topsep=0pt]
    \item Figure A1: $\alpha/\beta$ calibration iso-contour (M5 sweep, verifies H2a).
    % 【中文】图 A1：$\alpha/\beta$ 校准等值线（M5 扫描，验证 H2a）。
    \item Figure A2: $r_{\text{cp}}$ vs.\ Pass@1 scatter (M7 vs.\ M2, verifies H2b).
    % 【中文】图 A2：$r_{\text{cp}}$ vs.\ Pass@1 散点（M7 vs.\ M2，验证 H2b）。
    \item Figure A3: state\_match hit distribution, fail-streak vs.\ smooth (M6, verifies H2c).
    % 【中文】图 A3：state\_match 命中分布，失败连击 vs.\ 顺利（M6，验证 H2c）。
    \item Figure A4: M4 transfer leak-rate audit (abstract cognition vs.\ answer-fragment).
    % 【中文】图 A4：M4 迁移泄漏率审计（抽象认知 vs.\ 答案片段）。
\end{itemize}
All figures are \TBD{} pending evaluation.
% 【中文】所有图均为 \TBD{}，待评测。


\end{document}
